\chapter{Fluoroscopy}
\section*{What Is This Test?} Fluoroscopy uses continuous or pulsed X-rays to show moving anatomy and guide procedures in real time.
\section*{Alternative Names} Real-time X-ray imaging; fluoroscopic examination.
\section*{Principle} An X-ray beam and detector create live images; contrast may outline vessels, the digestive tract, joints, or urinary system.
\section*{Why This Test Is Ordered} It evaluates swallowing, gastrointestinal transit, joints, blood vessels, urinary structures, and guides injections, biopsies, or catheter procedures.
\section*{Primary vs Secondary Test} It is a targeted imaging study; ultrasound, CT, MRI, endoscopy, or plain radiography may be alternatives.
\section*{How This Test Is Performed} The patient is positioned and may drink contrast or receive contrast through a catheter while images are recorded.
\section*{What Preparation Is Needed} Preparation varies; fasting, bowel preparation, or medication instructions may apply. Report pregnancy possibility and contrast allergies.
\section*{Understanding Results} The report describes movement, leakage, narrowing, obstruction, abnormal anatomy, or procedure findings. Radiation dose depends on duration and procedure.
\section*{Follow-up Tests} Additional imaging, endoscopy, laboratory testing, or specialist consultation may follow.
\section*{References} \begin{enumerate}\item MedlinePlus. Fluoroscopy.\item U.S. Food and Drug Administration. Medical X-ray imaging information.\end{enumerate}
\clearpage\section*{Understanding Results and Follow-up}
The laboratory report should identify the specimen, method, units, reference interval or decision limit, and whether the result is positive, negative, abnormal, or inconclusive. Reference intervals vary with age, sex, pregnancy, specimen type, assay, and laboratory; a result outside a range is not automatically a diagnosis, and a result within a range does not always exclude disease. Discuss unexpected results with the ordering clinician, who can decide whether repeat or confirmatory testing, treatment, or referral is appropriate. Do not change medicines or delay urgent care based on an isolated result.
\section*{Regulatory Status and Authorization Date}This chapter describes a test category, not a specific commercial product. FDA, CDSCO, EU IVDR, and MHRA approval or registration dates therefore cannot be assigned without the assay manufacturer, product name, instrument, and intended use. For laboratory-developed tests, an individual agency approval date may not exist. See the Preface for the interpretation of regulatory status.
\clearpage
